\section{Skeleton entropy and input steering}\label{sec:skeleton}

The eighth paper excludes oblivious linear-time transposers whose work heads enter every cell boundedly often, given an
affirmative answer to Question~9.1 of the fifth paper, and extends the exclusion to adaptive machines of bounded entry
whose traces have entropy at most $\eta L$ for a fixed $\eta<\frac12$~\cite[Corollaries~6.3
and~6.4]{BoundedNesting2026}. It also observes that this entropy hypothesis is not robust: a preliminary sweep of the
input head with a one-step detour at every~$1$ of the first half of the input makes the trace reveal half the input
without changing the computation, and a form that survives such padding would bound the entropy of a coarse skeleton of
the trace instead \cite[Section~6.2]{BoundedNesting2026}. This section supplies that form, Theorem~SE, and draws its
consequences for heads that enter every cell boundedly often. At large scales the form is not affected by padding of
bounded entry before the run, such as that of the eighth paper; padding whose heads revisit cells unboundedly often can
still affect it (\cref{rem:skel-padding}).

The skeleton at scale~$w$ records the position of every head and the output count, rounded down to multiples of~$w$
(\cref{def:skeleton}). The tool is the block machine (\cref{def:block-machine,lem:block-machine}): on the inputs that
share a skeleton, the machine is an oblivious machine whose heads move on blocks of $w$ cells, its dependency DAG is the
walk DAG of the coarse walks, and bounded entry of a head gives bounded entry of its coarse walk. Theorem~F, the cone
lemma and the paired-distance argument of the fifth paper then apply to one recovery relation per skeleton class, which
gives Theorem~SE (\cref{thm:SE}). At large scales, heads of bounded entry carry almost no skeleton entropy
(\cref{lem:skeleton-count}). Hence, given [Q9.1], an affirmative answer to Question~9.1 for every degree bound, a
linear-time transposer whose work tapes have bounded entry must steer its input head by about half of the input at every
scale (Corollary~IS, \cref{cor:IS}), and no linear-time transposer, adaptive or not, has bounded entry on every tape,
the input tape included (Corollary~RV, \cref{cor:RV}). None of the results of this section is formalized; each has a
complete written proof. Only \cref{lem:skel-uniform}, Theorem~SE and Corollaries~IS and~RV assume [Q9.1].

\subsection{Skeletons}\label{sec:skel-skeletons}

We use the machine model of the eighth paper~\cite[Section~2.1]{BoundedNesting2026}. A machine~$M$ has a finite state
set~$\mathcal Q$, a read-only two-way input tape, $k\ge1$ two-way work tapes over a finite alphabet~$\Sigma$ that
contains the blank, and a write-only output tape that receives at most one bit per step; all heads start at cell~$0$ and
move by $-1$, $0$ or~$+1$, and every step writes the cell under each work head. At a length $L=d^2b(d)$ the input is a
word $X\in\{0,1\}^L$ with bits $X[p]$, $p\in[0,L)$, written in the cells $0,\dots,L-1$ of the input tape; every other
cell of the input tape reads as an end marker. (We write~$X$ for the input because $w$ will denote a scale.) On
input~$X$ the machine executes the steps $0,\dots,T(X)-1$ and then halts. For $0\le t\le T(X)$ we write $x_0(t)$ for the
position of the input head at time~$t$, that is, before step~$t$, $x_j(t)$ for the position of work head~$j$
($1\le j\le k$), and $n(t)$ for the number of bits appended at the steps before~$t$; we write $x_j(t;X)$ and $n(t;X)$
when the input varies. $M$ is a \emph{transposer} at length~$L$ if on every input it appends exactly $L$ bits and its
output bit at position~$r$ is $X[\tpos(r)]$, where $\tpos((di+j)b+s)=(dj+i)b+s$ for $0\le i,j<d$ and
$0\le s<b=b(d)$~\cite[Section~2.1]{BoundedNesting2026}; it works \emph{within $c(L+1)$ steps} if $T(X)\le c(L+1)$ for
every input. The map $\tpos$ is an involution; the fifth paper writes it~$\Pi$ \cite[(3.1)]{DirectedRecovery2026}.
Walks, hops, long hops, owner-load and walk DAGs are those of the eighth paper \cite[Sections~2.2
and~2.3]{BoundedNesting2026}, and entries are as in~\cite[Definition~2.1]{BoundedNesting2026}: an entry of a cell~$z$ by
a walk~$x$ is a time~$t$ with $x(t)=z$ and either $t=0$ or $x(t-1)\ne z$. The machine has \emph{entry at most~$V$ on
tape~$j$} on input~$X$ if the walk $x_j(\cdot\,;X)$ on $[0,T(X))$ enters every cell at most $V$ times. Entropies
$\mathrm H(\cdot)$ are Shannon entropies in bits, for a uniformly random input $X\in\{0,1\}^L$.

\begin{definition}[Skeletons]\label{def:skeleton}
Let $w\ge1$ be an integer, the \emph{scale}. The \emph{block} of $a\in\Z$ is $\mathcal B_a=\{wa,\dots,wa+w-1\}$, and on
the input tape $\mathcal B^{\mathrm{in}}_a=\mathcal B_a\cap[0,L)$. The \emph{coarse walk} of a walk~$x$ is
$\lfloor x/w\rfloor$; its value at~$t$ is the index of the block that contains $x(t)$. The \emph{skeleton} at scale~$w$
of the run on input~$X$ is the sequence
\[
  \sigma_w(X)=\Bigl(\bigl(\lfloor x_0(t)/w\rfloor,\lfloor x_1(t)/w\rfloor,\dots,\lfloor x_k(t)/w\rfloor,
  \lfloor n(t)/w\rfloor\bigr)\Bigr)_{0\le t<T(X)},
\]
and its \emph{input skeleton} is $\sigma^{\mathrm{in}}_w(X)=\bigl(\lfloor x_0(t)/w\rfloor\bigr)_{0\le t<T(X)}$; both
have length $T(X)$. The \emph{skeleton class} of a sequence~$\sigma$ is
$\mathcal X_\sigma=\{X\in\{0,1\}^L:\sigma_w(X)=\sigma\}$. For a skeleton~$\sigma$ we write $T_\sigma$ for its length and
$s_0,s_1,\dots,s_k,s_{\mathrm{out}}\colon[0,T_\sigma)\to\Z$ for its components; $s_1,\dots,s_k$ are the \emph{coarse
work walks} of~$\sigma$.
\end{definition}

All members of a class $\mathcal X_\sigma$ halt at time~$T_\sigma$, and at every time $t<T_\sigma$ their heads on one
tape lie in one block, as do their output counts; so their positions on each tape, and their output counts, differ by at
most $w-1$. Every component $s_j$ is a walk with $s_j(0)=0$, because
$\lfloor(z\pm1)/w\rfloor-\lfloor z/w\rfloor\in\{-1,0,1\}$, and $s_{\mathrm{out}}$ is nondecreasing with steps $0$
and~$1$. The trace of a run, that is, its head moves, append times and halting
time~\cite[Section~2.1]{BoundedNesting2026}, determines its skeleton at every scale; at $w=1$ the skeleton determines
the trace up to the moves of the last step and to whether that step appends. Finally, since
$\mathrm H(Y)\ge-\log_2\max_y\Pr[Y=y]$ for every random variable~$Y$ with finitely many values,
\begin{equation}\label{eq:skel-heavy}
  \mathrm H(\sigma_w(X))\le\bigl(\tfrac12-\eta\bigr)L\quad\Longrightarrow\quad
  \card{\mathcal X_\sigma}\ge2^{(\frac12+\eta)L}\ \text{for some skeleton }\sigma .
\end{equation}

\begin{lemma}[Coarse walks]\label{lem:skel-coarse}
\formal{MergeDepth/SkeletonWalk.lean}{MergeDepth.entries_skel_subset,MergeDepth.card_entries_skel_le,MergeDepth.card_block_changes_le}
Let $V,w,T\ge1$ be integers, let $x$ be a walk of length~$T$ with $x(0)=0$ that enters every cell at most $V$ times, and
let $s=\lfloor x/w\rfloor$.
\begin{enumerate}[label=(\alph*)]
\item If $t$ is an entry of an integer~$a$ by~$s$, then $x(t)\in\{wa,wa+w-1\}$ and $t$ is an entry of the cell $x(t)$
by~$x$.
\item The walk~$s$ enters every integer at most $2V$ times.
\item There are at most $2V\bigl((T-1)/w+2\bigr)-1$ times $t\in[1,T)$ with $s(t)\ne s(t-1)$.
\end{enumerate}
\end{lemma}

\begin{outline}
(a) If $t=0$, then $a=s(0)=0$ and $t$ is the entry of the cell $x(0)=0=wa$. If $t\ge1$, then $s(t-1)\ne a=s(t)$, so
$x(t-1)\notin\mathcal B_a\ni x(t)$. Hence $x(t)=x(t-1)\pm1$ is an end cell of $\mathcal B_a$, namely $wa$ or $wa+w-1$,
entered from outside the block, and $t$ is an entry of $x(t)$ by~$x$.

(b) By~(a), distinct entries of~$a$ by~$s$ are distinct entries by~$x$ of the at most two cells $wa$ and $wa+w-1$, each
of which is entered at most $V$ times.

(c) The cells visited by~$x$ form an interval of at most $T$ consecutive integers, and such an interval meets at most
$\lceil(T-1)/w\rceil+1\le(T-1)/w+2$ blocks. Every time $t\ge1$ with $s(t)\ne s(t-1)$ is an entry of $s(t)$ by~$s$.
By~(b) each of the visited blocks has at most $2V$ entries, and one of these entries, the entry of~$0$ at time~$0$, is
not such a time.

Formally, walks are maps $\N\to\Z$ with steps in $\{-1,0,1\}$, as in the eighth paper's library, the coarse walk is
\lean{MergeDepth.skel}, and the entries of a cell up to time~$T-1$ are the eighth paper's
\lean{BoundedNesting.entries}. In~(c) the formal bound uses the integer part of $(T-1)/w$, which is stronger.
\end{outline}

\begin{remark}[Hops of two members need not match]\label{rem:naive-skeleton}
A natural route to a robust form of Corollary~F$''$ applies Theorem~F to the trace of each member of a skeleton class
and merges the resulting cuts and cones. That route needs the hops of two members to correspond within $O(w)$ steps, and
this fails because a head can linger near a cell without leaving its block. Let $w\ge2$, $\lambda\ge2$ and $h\ge w$ be
integers and $a=2w+1$. Let $y$ be the walk that climbs from~$0$ to~$a$, stays at~$a$ for $\lambda$ steps, climbs to
$a+h$, which it reaches at the time $o=a+\lambda+h$, and descends to~$0$; for $a<s'<a+\lambda$ let $y_{s'}$ agree
with~$y$ except that $y_{s'}(s')=a-1$. The cells $a-1=2w$ and~$a$ lie in $\mathcal B_2$, so all these walks have the
same coarse walk; $y$ enters every cell at most twice and $y_{s'}$ at most three times. The upper long hops at the
cell~$2w$ that contain the time~$o$ start at the time $2w$ for~$y$ and at the time~$s'$ for~$y_{s'}$, so their starts
differ by $s'-2w$, which takes every value from~$2$ to~$\lambda$, although the cell $2w$ is farther than $w$ from both
ends of the family of~$y$ with owner~$o$, whose base is~$0$ and whose height is $a+h$. Pauses are not essential when
$w\ge3$. For $h>\lambda$ let $y'$ climb from~$0$ to $2w+2$, make an excursion of $2\lambda$ steps up to $2w+2+\lambda$
and back, visit $2w+1,2w+2,2w+1,2w+2$, then climb to $2w+2+h$ and descend to~$0$, and let $y''$ agree with~$y'$ except
that it visits $2w+1,2w,2w+1,2w+2$ after the excursion. The cells $2w,2w+1,2w+2$ lie in $\mathcal B_2$ because $w\ge3$,
so $y'$ and~$y''$ have the same coarse walk; neither pauses, both enter every cell at most five times, and their upper
long hops at the cell~$2w$ that contain the time of the top start $2\lambda+4$ steps apart. The block machine below
avoids per-trace structure altogether: its cut is determined by the skeleton.
\end{remark}

\subsection{The block machine}\label{sec:skel-block}

Fix a scale~$w$ and a skeleton~$\sigma$ with $\mathcal X_\sigma\ne\emptyset$, and write $T=T_\sigma$.

\begin{definition}[The block machine]\label{def:block-machine}
For $X\in\mathcal X_\sigma$ and $0\le t<T$, the \emph{block value} $\bar v_t(X)$ consists of
\begin{enumerate}[label=(\roman*)]
\item the state after step~$t$;
\item for every work tape~$j$, the contents of the $w$ cells of $\mathcal B_{s_j(t)}$ after step~$t$;
\item for every head $j\in\{0,1,\dots,k\}$, the offset $x_j(t+1)-ws_j(t)$;
\item the word $O_t(X)$ of the output bits at the positions in $[ws_{\mathrm{out}}(t),n(t+1))$, in order.
\end{enumerate}
For $1\le j\le k$ and $0\le t<T$ let $\prev^\sigma_j(t)=\max\{t'<t:s_j(t')=s_j(t)\}$ when this set is nonempty. The
\emph{block DAG} $\bar H_\sigma$ has vertex set $[0,T)$, a \emph{state edge} $t\to t-1$ for every $1\le t<T$, and a
\emph{block edge} $t\to\prev^\sigma_j(t)$ for every work tape~$j$ and every~$t$ for which $\prev^\sigma_j(t)$ is
defined; an edge is determined by its endpoints, so a block edge of length one is a state edge. The block values on
$\mathcal X_\sigma$, the block DAG and the components of~$\sigma$ form the \emph{block machine} of~$M$ on~$\sigma$. Its
\emph{value width} is
\[
  \beta_w=\clog{\card{\mathcal Q}}+kw\,\clog{\card\Sigma}+(k+1)\clog{(w+2)}+w+1 .
\]
\end{definition}

The offset in~(iii) lies in $\{-1,0,\dots,w\}$, because $x_j(t)\in\mathcal B_{s_j(t)}$ and heads move by at most one
cell, and the word in~(iv) has length at most~$w$, because $n(t+1)\le n(t)+1\le ws_{\mathrm{out}}(t)+w$. So a block
value is described by at most $\beta_w$ bits: (ii) takes $kw\clog{\card\Sigma}$ bits, (iii) takes $(k+1)\clog{(w+2)}$
bits, and (iv), one of the $2^{w+1}-1$ binary words of length at most~$w$, takes $w+1$ bits. Note that $\beta_w\ge6$.
Genies, $\Gamma$-free paths and full depth are defined for $\bar H_\sigma$ as for the dependency DAG of a trace in the
eighth paper \cite[Section~2.4]{BoundedNesting2026}: a \emph{genie} is a set $\Gamma\subseteq[0,T)$, it reveals the
block values $\bar v_u(X)$, $u\in\Gamma$, with $\beta_w$ bits each, at \emph{cost} $\beta_w\card\Gamma$; a path
$t_0\to\dots\to t_m$ is \emph{$\Gamma$-free} if $t_1,\dots,t_m\notin\Gamma$; and $\Gamma$ has \emph{full depth at
most~$D$} if every $\Gamma$-free path has at most $D$ edges.

\begin{lemma}[The block machine]\label{lem:block-machine}
\notformalized
Let $\sigma$ be a skeleton at scale~$w$ with $\mathcal X_\sigma\ne\emptyset$, and $T=T_\sigma$.
\begin{enumerate}[label=(\alph*)]
\item \emph{(Dependency DAG.)} $\bar H_\sigma$ is the walk DAG of the coarse work walks $s_1,\dots,s_k$, and every
vertex has in-degree and out-degree at most $k+1$.
\item \emph{(Obliviousness.)} For every $t<T$ there is a function~$\Phi_t$, determined by $M$, $\sigma$ and~$t$, such
that for every $X\in\mathcal X_\sigma$
\[
  \bar v_t(X)=\Phi_t\Bigl(\bar v_{t-1}(X),\ \bar v_{\prev^\sigma_1(t)}(X),\dots,\bar v_{\prev^\sigma_k(t)}(X),\
  X|_{\mathcal B^{\mathrm{in}}_{s_0(t)}}\Bigr),
\]
where an argument that is not defined, because $t=0$ or because $\prev^\sigma_j(t)$ is not defined, is replaced by a
fixed symbol.
\item \emph{(Output decoding.)} Suppose that every member of $\mathcal X_\sigma$ appends exactly $L$ bits. For
$r\in[0,L)$ the time $\hat t_r=\max\{t<T:s_{\mathrm{out}}(t)=\lfloor r/w\rfloor\}$ exists, and for every
$X\in\mathcal X_\sigma$ the output bit at position~$r$ is the letter of $O_{\hat t_r}(X)$ at place
$r-w\lfloor r/w\rfloor$, counted from~$0$. At most $w$ positions~$r$ share one time~$\hat t_r$.
\item \emph{(Entries.)} If some member of $\mathcal X_\sigma$ has entry at most~$V$ on every work tape, then every
coarse work walk~$s_j$ enters every integer at most $2V$ times, and the set of all long hops of $s_1,\dots,s_k$ has
owner-load at most $4V-1$~\cite[Definition~2.2]{BoundedNesting2026}.
\end{enumerate}
\end{lemma}

\begin{informalproof}
(a) For a walk~$y$ and a time~$u$, the hop of~$y$ ending at~$u$ is $(\prev_y(u),u)$, where
$\prev_y(u)=\max\{t<u:y(t)=y(u)\}$, and it is long if $u-\prev_y(u)\ge2$~\cite[Section~2.2]{BoundedNesting2026}. For
$y=s_j$ this is $\prev^\sigma_j$. So the block edges of length at least two are the edges $u\to p$ for the long hops
$(p,u)$ of $s_1,\dots,s_k$, which are the hop edges of their walk DAG, and the block edges of length one are pauses,
which the walk DAG identifies with state edges~\cite[Section~2.3]{BoundedNesting2026}, as $\bar H_\sigma$ does. So the
two DAGs coincide. A vertex~$t$ has out-edges only to $t-1$ and to the at most $k$ vertices $\prev^\sigma_j(t)$. Its
in-edges come from $t+1$ and, for each~$j$, from at most one~$t'$ with $\prev^\sigma_j(t')=t$: if $t<t'<t''$ and
$\prev^\sigma_j(t')=\prev^\sigma_j(t'')=t$, then $s_j(t'')=s_j(t)=s_j(t')$, so $\prev^\sigma_j(t'')\ge t'>t$.

(b) Fix $t<T$ and $X\in\mathcal X_\sigma$. Step~$t$ applies the transition function of~$M$ to the state at time~$t$, the
input symbol under $x_0(t)$ and the symbols under $x_1(t),\dots,x_k(t)$. These data, and the positions $x_j(t)$, are
functions of the arguments of~$\Phi_t$ through maps that depend only on $M$, $\sigma$ and~$t$:
\begin{itemize}
\item The state at time~$t$ is the initial state if $t=0$, and item~(i) of $\bar v_{t-1}(X)$ otherwise.
\item $x_j(0)=0$, and for $t\ge1$, $x_j(t)=ws_j(t-1)+e_j$, where $e_j$ is the offset of head~$j$ in item~(iii) of
$\bar v_{t-1}(X)$ and $s_j(t-1)$ is read from~$\sigma$.
\item Let $1\le j\le k$. If $t'=\prev^\sigma_j(t)$ is defined, the contents of $\mathcal B_{s_j(t)}$ on tape~$j$ at
time~$t$ are those after step~$t'$, which form item~(ii) of $\bar v_{t'}(X)$: for $t'<\tau<t$ we have
$s_j(\tau)\ne s_j(t)$, so $x_j(\tau)\notin\mathcal B_{s_j(t)}$, and tape~$j$ is written only by head~$j$ and only in the
cell under it. If $\prev^\sigma_j(t)$ is not defined, head~$j$ has not been in $\mathcal B_{s_j(t)}$ before time~$t$, so
the block is blank. In both cases the symbol under head~$j$ is the one at place $x_j(t)-ws_j(t)$ of the block.
\item If $x_0(t)\in[0,L)$, then $x_0(t)\in\mathcal B^{\mathrm{in}}_{s_0(t)}$, and the input symbol $X[x_0(t)]$ is a
component of $X|_{\mathcal B^{\mathrm{in}}_{s_0(t)}}$; otherwise it is the end marker.
\end{itemize}
The transition yields the new state, the symbols written, the moves and the appended bit, if any. Item~(i) of
$\bar v_t(X)$ is the new state; item~(ii) for tape~$j$ is the contents of $\mathcal B_{s_j(t)}$ at time~$t$ with the
cell $x_j(t)$ overwritten by the symbol written there; item~(iii) for head~$j$ is $x_j(t)-ws_j(t)$ plus the move of
head~$j$. For item~(iv): if $t\ge1$ and $s_{\mathrm{out}}(t)=s_{\mathrm{out}}(t-1)$, then $O_{t-1}(X)$ is the word of
the positions in $[ws_{\mathrm{out}}(t),n(t))$, and $O_t(X)$ is $O_{t-1}(X)$ followed by the bit appended at step~$t$,
if any, whose position is $n(t)$. Otherwise $O_t(X)$ consists of that bit alone, or is empty, because
$n(t)=ws_{\mathrm{out}}(t)$: for $t=0$ both sides are~$0$, and for $t\ge1$ the count grows by at most one per step, so
$s_{\mathrm{out}}(t)\ne s_{\mathrm{out}}(t-1)$ gives $n(t-1)<ws_{\mathrm{out}}(t)\le n(t)\le n(t-1)+1$. The values of
$s_{\mathrm{out}}$ are read from~$\sigma$.

(c) Fix $X\in\mathcal X_\sigma$. Since $n(0)=0$, $n(T)=L$ and $n$ grows by $0$ or~$1$ per step, for every $r<L$ there is
a step $t_r<T$ with $n(t_r)=r<n(t_r+1)$, the step that appends bit~$r$. Then $s_{\mathrm{out}}(t_r)=\lfloor r/w\rfloor$,
so $\hat t_r$ exists and $\hat t_r\ge t_r$, whence $n(\hat t_r+1)\ge n(t_r+1)=r+1$. So $O_{\hat t_r}(X)$, the word of
the positions in $[w\lfloor r/w\rfloor,n(\hat t_r+1))$, contains position~$r$, at place $r-w\lfloor r/w\rfloor$. The
time $\hat t_r$ depends only on $\sigma$ and~$r$, and $s_{\mathrm{out}}(\hat t_r)=\lfloor r/w\rfloor$; so positions with
distinct values of $\lfloor r/w\rfloor$ have distinct times, and at most $w$ positions share a value.

(d) The coarse work walks of~$\sigma$ are those of the head walks $x_j(\cdot\,;X)$ on $[0,T)$ of that member~$X$, which
start at~$0$. By \cref{lem:skel-coarse}(b) they enter every integer at most $2V$ times, and the eighth paper's
Corollary~3.2, which is Lean-verified there, applied with $2V$ in place of~$V$, bounds the owner-load of the set of all
their long hops by $2\cdot2V-1$~\cite[Corollary~3.2]{BoundedNesting2026}.
\end{informalproof}

So on a skeleton class, $M$ is an oblivious machine of value width~$\beta_w$: its work heads move along the coarse work
walks on tapes whose cells are blocks, its input head reads one block of $w$ input cells per step, its output is written
one block at a time, and its dependency DAG is $\bar H_\sigma$. Motion inside a block, however data-dependent, creates
no edge: it is carried by the items~(ii) and~(iii) of the block values along state edges. This is what makes adaptivity
tractable: the tools for oblivious machines apply to each skeleton class, at the price of the value width~$\beta_w$, and
only the coarse walks enter the depth analysis.

\subsection{One recovery relation per skeleton class}\label{sec:skel-class}

\begin{lemma}[A genie for a skeleton class]\label{lem:skel-genie}
\notformalized
Let $M$ be a machine with $k$ work tapes, let $V,w\ge1$ be integers, $c>0$ and $0<\eps\le1$, and put $C=\max(c,1)$,
\[
  \eps''=\frac{\eps}{4\beta_wC},\qquad P=\Bigl\lceil\frac{4\beta_wC}{\eps}\Bigr\rceil,\qquad
  D_1=\DF(k,4V-1,\eps''),\qquad D=D_1P-1,
\]
where $\DF$ is the bound of Theorem~F~\cite[Theorem~3.6]{BoundedNesting2026}. Let $\sigma$ be a skeleton of~$M$ at
scale~$w$ at a length $L\ge1$, with $T_\sigma\le c(L+1)$, and suppose that some member of $\mathcal X_\sigma$ has entry
at most~$V$ on every work tape. Then there is a genie $\Gamma\subseteq[0,T_\sigma)$, depending only on~$\sigma$, of cost
$\beta_w\card\Gamma\le\eps L$ and full depth at most~$D$ in~$\bar H_\sigma$.
\end{lemma}

\begin{informalproof}
Write $T=T_\sigma$. Since $\beta_w\ge6$, $\eps''\in(0,1]$. By \cref{lem:block-machine}(d) the set of all long hops of
the coarse work walks $s_1,\dots,s_k$, which are walks of length~$T$, has owner-load at most $4V-1$. Theorem~F
\cite[Theorem~3.6]{BoundedNesting2026}, which is Lean-verified there and allows pauses and any coupling of the walks,
gives a cut~$F$ of their walk DAG with $\card F\le\eps''T$ such that every path that uses no edge of~$F$ has at most
$D_1-1$ hop edges. By \cref{lem:block-machine}(a) this walk DAG is $\bar H_\sigma$, and its hop edges are its edges of
length at least two. Fix such a cut; it depends only on~$\sigma$.

Let $\Gamma_0$ be the set of the vertices~$u$ such that some edge $t\to u$ lies in~$F$; so $\card{\Gamma_0}\le\card F$.
Every edge of a $\Gamma_0$-free path ends outside~$\Gamma_0$, so it is not in~$F$, and such a path has at most $D_1-1$
edges of length at least two. Put $\Gamma=\Gamma_0\cup\{t<T:t\equiv P-1\pmod P\}$ and let $\pi$ be a $\Gamma$-free path.
Removing its at most $D_1-1$ edges of length at least two leaves at most $D_1$ maximal runs of state edges. The vertices
of such a run other than its first are vertices of~$\pi$ other than its first, hence lie outside~$\Gamma$, and they are
consecutive integers; so there are at most $P-1$ of them, and the run has at most $P-1$ edges. Hence $\pi$ has at most
$(D_1-1)+D_1(P-1)=D$ edges. This is the long-edge conversion of the fourth paper~\cite[Lemma~4.21]{FourthPaper2026} with
$\mu=1$, whose proof uses only the DAG. Finally $\card\Gamma\le\card F+\lfloor T/P\rfloor$, so, using $T\le C(L+1)$ and
$L\ge1$,
\[
  \beta_w\card\Gamma\le\beta_w\eps''T+\frac{\beta_wT}P\le\frac{\eps T}{4C}+\frac{\eps T}{4C}\le\frac{\eps(L+1)}2\le\eps L.
  \qedhere
\]
\end{informalproof}

This is the eighth paper's conversion of cuts into genies~\cite[Lemma~2.6]{BoundedNesting2026}, run on the block
machine. The genie is the same set of times for every member of the class; only the revealed values depend on the input.

\begin{lemma}[Cones on a skeleton class]\label{lem:skel-cones}
\notformalized
Let $\sigma$ be a skeleton at scale~$w$ with $\mathcal X_\sigma\neq\emptyset$, every member of which appends exactly $L$
bits, and let $\Gamma\subseteq[0,T_\sigma)$ have full depth at most~$D$ in~$\bar H_\sigma$. For $t<T_\sigma$ let
$C_\Gamma(t)$ be the set of endpoints of $\Gamma$-free paths from~$t$ in~$\bar H_\sigma$, so that $t\in C_\Gamma(t)$,
and for $r,q\in[0,L)$ put
\[
  J_\Gamma(t)=\bigcup_{u\in C_\Gamma(t)}\mathcal B^{\mathrm{in}}_{s_0(u)},\qquad I_\Gamma(r)=J_\Gamma(\hat t_r),\qquad
  N_\Gamma(q)=I_\Gamma(\tpos(q)).
\]
Let $G_\Gamma(X)=(\bar v_u(X))_{u\in\Gamma}$ be the revealed word, and $B=\sum_{h=0}^D(k+1)^h$.
\begin{enumerate}[label=(\alph*)]
\item $\card{C_\Gamma(t)}\le B$ for every~$t$, and $\card{N_\Gamma(q)}\le wB$ for every~$q$.
\item If $X,Y\in\mathcal X_\sigma$ satisfy $G_\Gamma(X)=G_\Gamma(Y)$ and $X|_{J_\Gamma(t)}=Y|_{J_\Gamma(t)}$, then
$\bar v_t(X)=\bar v_t(Y)$. Consequently, if $G_\Gamma(X)=G_\Gamma(Y)$ and $X|_{I_\Gamma(r)}=Y|_{I_\Gamma(r)}$, the
output bits of $X$ and~$Y$ at position~$r$ agree.
\item If $M$ transposes every input of $\mathcal X_\sigma$, then for every value~$g$ of the revealed word the fibre
$\{X\in\mathcal X_\sigma:G_\Gamma(X)=g\}$ is a binary storage code for the recovery relation $N_\Gamma$ on $[0,L)$ in
the sense of the fifth paper~\cite[Definition~2.1]{DirectedRecovery2026}.
\end{enumerate}
\end{lemma}

\begin{informalproof}
(a) Every vertex of $C_\Gamma(t)$ is the endpoint of a $\Gamma$-free path from~$t$ with at most $D$ edges, and by
\cref{lem:block-machine}(a) at most $(k+1)^h$ paths with $h$ edges start at~$t$. Every block has at most $w$ input
cells, and $N_\Gamma(q)=J_\Gamma(\hat t_{\tpos(q)})$.

(b) By strong induction on~$t$, as in the fourth paper's cone lemma~\cite[Lemma~4.4]{FourthPaper2026}, with
\cref{lem:block-machine}(b) in place of its local recursion. We show that the arguments of~$\Phi_t$ agree on $X$
and~$Y$. The constant arguments agree. Since $t\in C_\Gamma(t)$, the block $\mathcal B^{\mathrm{in}}_{s_0(t)}$ is
contained in $J_\Gamma(t)$, on which $X$ and~$Y$ agree. Every other argument is $\bar v_{t'}$ for an edge $t\to t'$ of
$\bar H_\sigma$. If $t'\in\Gamma$, then $\bar v_{t'}(X)=\bar v_{t'}(Y)$ is a component of the common revealed word. If
$t'\notin\Gamma$, then every $\Gamma$-free path from~$t'$, preceded by the edge $t\to t'$, is a $\Gamma$-free path
from~$t$; so $C_\Gamma(t')\subseteq C_\Gamma(t)$ and $J_\Gamma(t')\subseteq J_\Gamma(t)$, and the induction hypothesis
for $t'<t$ gives $\bar v_{t'}(X)=\bar v_{t'}(Y)$. Hence $\bar v_t(X)=\bar v_t(Y)$. The second sentence is the case
$t=\hat t_r$, by \cref{lem:block-machine}(c).

(c) Let $X$ and~$Y$ lie in the fibre and agree on $N_\Gamma(q)=I_\Gamma(\tpos(q))$. By~(b) their output bits at position
$\tpos(q)$ agree; since $M$ transposes and $\tpos$ is an involution, these bits are $X[q]$ and~$Y[q]$.
\end{informalproof}

The relation $N_\Gamma$ depends only on $\sigma$ and~$\Gamma$. Compared with the cone relation of a single trace, every
cone vertex contributes a block of $w$ input cells instead of one, and the at most $w$ output positions of one output
block share their cone.

\subsection{Paired distances on a skeleton class}\label{sec:skel-paired}

The fifth paper proves that, after the paid deletion of the coordinates of high reuse, the cone relations of every
linear-time schedule with a genie of fixed full depth satisfy the paired-distance condition \cite[Theorem~3.3 and
Corollaries~3.7 and~3.8]{DirectedRecovery2026}, and that the argument applies to a single frozen
trace~\cite[Remark~3.9]{DirectedRecovery2026}. The relation $N_\Gamma$ of a skeleton class is not the cone relation of
one trace, so we prove the three steps of that argument for it. There are three changes: the relaxation carries $2w-1$
heads on each tape around the heads of one member of the class, the occupancy bound becomes $wB+1$ instead of $B+1$, and
the initial potential becomes $m\log_23$ instead of~$m$.

We use the geometric relaxation of the fifth paper~\cite[Section~3.2]{DirectedRecovery2026}, with several heads allowed
on one tape. There are finitely many tapes with cells indexed by~$\Z$, among them an \emph{input tape} and an
\emph{output tape}; $K$ unit-speed heads, distributed over the tapes in any way; and $\rho$ registers, which have no
addresses. The \emph{slots} are the cells and the registers. A \emph{tick} first transfers labels arbitrarily among the
cells under heads and the registers, leaving all other slots unchanged, and then moves every head by at most one cell;
moving a head does not move labels. Every label follows one trajectory through the slots; trajectories may share a slot
and separate later. The \emph{occupancy} of a slot is the number of labels in it, and it is bounded at every
\emph{placement}, that is, before and after the transfers of every tick. The fifth paper uses one head per tape, all
starting at cell~$0$.

\begin{lemma}[Transport with several heads per tape]\label{lem:skel-transport}
\notformalized
Let $d\ge2$ and $b\ge1$ be integers, $L=d^2b$ and $W=db$, let $\omega\ge0$ be an integer with $2\omega+1\le W$, and let
$S\subseteq[0,L)$ have $m>0$ elements. In the relaxation with $K$ heads and $\rho$ registers, suppose that trajectories
start at the distinct input cells $p\in S$, with no labels elsewhere, and end at the distinct output cells $\tpos(p)$;
that every head starts in $[-\omega,\omega]$ on its tape; that every slot has occupancy at most~$\kappa$ at every
placement; and that $U$ ticks are used. Then
\[
  m\log_2d\le\kappa K(K+\rho+3)U+3L .
\]
\end{lemma}

\begin{informalproof}
We follow the proof of the fifth paper's Lemma~3.4~\cite{DirectedRecovery2026}. Fix the trajectories. For
$\theta\in\{0,\dots,W-1\}$ partition every tape into the \emph{bins} $\{z:\lfloor(z-\theta)/W\rfloor=i\}$, $i\in\Z$, the
tape being part of the name of a bin. At every placement the \emph{active class} is the union of the registers and of
the bins that contain at least one head, and every other bin is a class of its own. The active class has at most
$KW+\rho$ slots, hence at most $\kappa(KW+\rho)$ labels. Let $p^*$ be uniform in~$S$, let $A=\lfloor p^*/W\rfloor$ be
its input row, let $Z_\theta$ be the class that contains the label of~$p^*$ at a given placement, and put
$\Phi=m\,\mathbb E_\theta\,\mathrm H(A\mid Z_\theta)$. For random variables $Z,Z'$ with finitely many values,
$\mathrm H(A\mid Z')-\mathrm H(A\mid Z)\le\mathrm H(Z\mid Z')$, because
$\mathrm H(A\mid Z')\le\mathrm H(A,Z\mid Z')\le\mathrm H(Z\mid Z')+\mathrm H(A\mid Z)$
\cite[(3.6)]{DirectedRecovery2026}.

\emph{Transfers.} Both ends of every transfer are head cells or registers, hence in the active class, so no label
changes class and $\Phi$ does not change.

\emph{Moves.} Process the moves of a tick one head at a time. A unit move of a head from~$z$ to $z\pm1$ changes the bin
of that head for exactly one~$\theta$; for every other~$\theta$ the classes are unchanged. For that~$\theta$, the
entered bin joins the active class if it is not already in it, the left bin leaves it if no other head is in it, and all
other classes are unchanged. So every label of the new active class had one of two old class names, the old active class
or the entered bin, and every label of a released bin had the old active name. Hence
$\mathrm H(Z_{\mathrm{old}}\mid Z_{\mathrm{new}})$ is at most the probability that $Z_{\mathrm{new}}$ is the active
class, which is at most $\kappa(KW+\rho)/m$, and
$m\,\mathrm H(A\mid Z_{\mathrm{new}})-m\,\mathrm H(A\mid Z_{\mathrm{old}})\le\kappa(KW+\rho)$. Averaged over~$\theta$,
the move raises $\Phi$ by at most $\kappa(KW+\rho)/W\le\kappa(K+\rho)$. A tick has at most $K$ unit moves, so over the
$U$ ticks $\Phi$ grows by at most $\kappa K(K+\rho)U$.

\emph{Start.} All labels lie in their input cells. The heads of the input tape lie in $[-\omega,\omega]$, at most $W$
consecutive cells, which meet at most two adjacent bins; their union, $2W$ consecutive cells, meets at most three of the
rows $\{iW,\dots,iW+W-1\}$. A bin that is not active meets at most two rows, and the active bins of other tapes and the
registers hold no labels. So for every~$\theta$, given the class, $A$ takes at most three values, and initially
$\Phi\le m\log_23$.

\emph{End.} The label of~$p$ sits in output cell $\tpos(p)$. For $p=(di+j)b+s$ the row $\lfloor p/W\rfloor$ is~$i$,
and $\tpos(p)=(dj+i)b+s$ gives $i=\lfloor(\tpos(p)\bmod W)/b\rfloor$: the row is a $W$-periodic function of the output cell
that takes every value on exactly $b$ cells of a period. So a bin of the output tape that holds $n'$ labels has row
multiplicities $n'_a\le b$, and
$n'\,\mathrm H(A\mid\text{this bin})=n'\log_2n'-\sum_an'_a\log_2n'_a\ge n'\log_2n'-n'\log_2b$. At most $L/W+2$ bins meet
$[0,L)$, so by the convexity of $y\log_2y$, if every bin is a class of its own, then
$m\,\mathrm H(A\mid Z_\theta)\ge m\log_2\bigl(m/(b(L/W+2))\bigr)=m\log_2d-m\log_2\bigl((L+2W)/m\bigr)$. The actual
partition merges the active bins and the registers into one class; this coarsening cannot lower the conditional entropy.
So finally $\Phi\ge m\log_2d-m\log_2((L+2W)/m)$.

\emph{Combination.} Hence $m\log_2d\le\kappa K(K+\rho)U+m\log_23+m\log_2((L+2W)/m)$. Since $d\ge2$, we have
$L+2W\le2L$, and $y\log_2(a/y)\le a/(e\ln2)$ for $0<y\le a$; so $m\log_2((L+2W)/m)\le2L/(e\ln2)<1.07L$, while
$m\log_23<1.59L$ because $m\le L$. This gives the claim.
\end{informalproof}

\begin{lemma}[Realizing a matching of incidences]\label{lem:skel-matching}
\notformalized
Let $\sigma$, $\Gamma$, $D$, $B$ and $N_\Gamma$ be as in \cref{lem:skel-cones}, write $T=T_\sigma$, and let $E$ be the
incidence graph of $N_\Gamma$: it has a vertex~$L_p$ and a vertex~$R_p$ for every $p\in[0,L)$, and an edge $L_pR_q$
exactly when $p\in N_\Gamma(q)$~\cite[Definitions~2.3 and~3.2]{DirectedRecovery2026}. Fix $X^*\in\mathcal X_\sigma$,
with head walks $z_j(t)=x_j(t;X^*)$ for $0\le j\le k$ and output count $z_{\mathrm{out}}(t)=n(t;X^*)$, $0\le t\le T$.
Consider the relaxation with an input tape, $k$ work tapes and an output tape, each carrying $2w-1$ heads, which at
tick~$t$ occupy the cells $z(t)+i$, $\abs i\le w-1$, where $z$ is the tape's reference walk $z_0$, $z_j$ or
$z_{\mathrm{out}}$; and with one register. So $K=K_w:=(2w-1)(k+2)$, $\rho=1$, every head starts in $[-(w-1),w-1]$, and
the head configurations at all ticks are fixed in advance. Let $\mathcal M$ be a matching in~$E$. Then there is a
realization in $T$ ticks that carries, for every edge $L_pR_q\in\mathcal M$, one label from input cell~$p$ to output
cell $\tpos(q)$, with occupancy at most $wB+1$. It can be reversed in $2T$ ticks, from the head configuration after the
last tick back to the initial one, with the same occupancy.
\end{lemma}

\begin{informalproof}
We follow the proof of the fifth paper's Lemma~3.5~\cite{DirectedRecovery2026}.

\emph{Coverage.} Since $X^*\in\mathcal X_\sigma$, we have $z_j(t)\in\mathcal B_{s_j(t)}$ for $0\le j\le k$ and
$z_{\mathrm{out}}(t)\in\mathcal B_{s_{\mathrm{out}}(t)}$ for $t<T$. The $2w-1$ cells $z(t)+i$, $\abs i\le w-1$, around a
cell of a block contain the whole block. So at tick~$t$ every cell of $\mathcal B_{s_j(t)}$ on tape~$j$, for
$0\le j\le k$, and every cell of $\mathcal B_{s_{\mathrm{out}}(t)}$ on the output tape lies under a head.

\emph{Paths and slots.} For every edge $L_pR_q\in\mathcal M$ we have $p\in N_\Gamma(q)=J_\Gamma(\hat t_{\tpos(q)})$;
choose a $\Gamma$-free path $\hat t_{\tpos(q)}=t_0\to t_1\to\dots\to t_\ell=u$ in $\bar H_\sigma$ with
$p\in\mathcal B^{\mathrm{in}}_{s_0(u)}$, so $\ell\le D$. Assign a slot to every edge of $\bar H_\sigma$: to a state edge
the register, and to an edge $t\to t'$ of length at least two the \emph{base cell} $ws_j(t)=ws_j(t')$ of tape~$j$, for
the least~$j$ with $t'=\prev^\sigma_j(t)$. By coverage, the slot of an edge $t\to t'$ is the register or a head cell at
the ticks $t'$ and~$t$.

\emph{The realization.} Each label moves along its path read forwards in time. The label of $L_pR_q$ stays in input
cell~$p$ until tick~$u$. At tick~$t_i$, before the moves, it is transferred from the slot by which it arrives (input
cell~$p$ if $i=\ell$, and the slot of the edge $t_i\to t_{i+1}$ otherwise) to the slot by which it leaves (the slot of the
edge $t_{i-1}\to t_i$ if $i\ge1$, and output cell $\tpos(q)$ if $i=0$). All these slots are the register or head cells
at tick~$t_i$: input cell~$p$ lies in $\mathcal B^{\mathrm{in}}_{s_0(u)}$; output cell $\tpos(q)$ lies in the block
$\mathcal B_{s_{\mathrm{out}}(\hat t_{\tpos(q)})}$ by the definition of $\hat t_{\tpos(q)}$; and edge slots are covered
as said. Between the ticks $t'$ and~$t$ of an edge of length at least two, its labels wait in its base cell $ws_j(t)$.
That cell may lie under a head at the ticks strictly between, because the heads reach into neighbouring blocks, but no
transfer touches it then: a transfer at tick~$\tau$ uses, on tape~$j$, only the base cell $ws_j(\tau)$, and
$s_j(\tau)\ne s_j(t)$ for $t'<\tau<t$. All transfers of a tick precede its moves, so $T$ ticks suffice.

\emph{Occupancy.} A chosen path through a vertex~$t$ starts at a time $\hat t_r$ from which $t$ can be reached by a path
with at most $D$ edges. Since in-degrees are at most $k+1$ (\cref{lem:block-machine}(a)), there are at most $B$ such
times; at most $w$ positions~$r$ share each of them (\cref{lem:block-machine}(c)); and every~$R_q$, hence every
$r=\tpos(q)$, lies on at most one edge of~$\mathcal M$. So at most $wB$ chosen paths pass through~$t$. Before the
transfers of tick~$\tau$ the register holds only labels of the edge $\tau\to\tau-1$, and after them only labels of the
edge $\tau+1\to\tau$; all these paths pass through one vertex, so there are at most $wB$ such labels. Before the
transfers of tick~$\tau$, the base cell of block~$a$ on tape~$j$ holds only labels of an edge of tape~$j$ that comes
from a long hop of $s_j$ at~$a$ spanning~$\tau$, and after them only labels of such an edge from a hop spanning
$\tau+1$; by property~(P1) of the walk~$s_j$~\cite[Section~2.2]{BoundedNesting2026}, at most one hop of $s_j$ at~$a$
spans a given time, so these labels belong to paths through one vertex, and there are at most $wB$ of them. An input
cell holds at most its own label, and an output cell at most the label that ends there, because $\mathcal M$ is a
matching. So the occupancy is at most $wB+1$ at every placement.

\emph{Reversal.} The heads follow the walks $z+i$, whatever $\mathcal M$ is. A tick is reversed in two ticks: one undoes
its moves with all labels held, and one undoes its transfers with the heads fixed; the ends of each reversed transfer
are again head cells or the register. The reversal passes through the placements of the realization in the opposite
order, so it has the same occupancy.
\end{informalproof}

\begin{lemma}[Disjoint short paired paths]\label{lem:skel-disjoint}
\notformalized
In the setting of \cref{lem:skel-matching}, suppose that $L=d^2b$ with $d\ge2$ and $b=b(d)$, that $db\ge2w-1$, and that
$T\le c(L+1)$. Let $s\ge1$ be odd. If $E$ contains $m$ pairwise vertex-disjoint paths with $s$ edges each, every one of
which joins $L_p$ to~$R_p$ for some~$p$, then
\[
  m\le A^{(w)}_s\frac L{\log_2d},\qquad A^{(w)}_s=4sc(wB+1)K_w(K_w+4)+3 .
\]
\end{lemma}

\begin{informalproof}
We follow the proof of the fifth paper's Lemma~3.6~\cite{DirectedRecovery2026}. Write $s=2a+1$, and list the vertices of
each path from its left end: $L_{y_0},R_{y_1},L_{y_2},\dots,R_{y_s}$, with $y_0=y_s=p$. Since the paths are
vertex-disjoint, for each $i<s$ their $i$th edges form a matching $\mathcal M_i$ of~$E$, consisting of the edges
$L_{y_i}R_{y_{i+1}}$ if $i$ is even and $L_{y_{i+1}}R_{y_i}$ if $i$ is odd. Run, in this order, the realization of
\cref{lem:skel-matching} for $\mathcal M_0$, the reversed realization for $\mathcal M_1$, the realization for
$\mathcal M_2$, and so on, ending with the realization for $\mathcal M_{2a}$. For even~$i$ the label of a path moves
from input cell $y_i$ to output cell $\tpos(y_{i+1})$, and for odd~$i$ the reversed realization moves it from output
cell $\tpos(y_i)$ to input cell $y_{i+1}$. So consecutive stages meet at the same physical cell, and $\tpos$ is never
applied as a free operation. The forward stages run from the initial head configuration to the final one and the
reversed stages back, so the configurations agree at every junction, and at every junction the labels occupy distinct
cells, so the occupancy stays at most $wB+1$. The result carries the label of each~$p$ from input cell~$p$ to output
cell $\tpos(p)$, with every head starting in $[-(w-1),w-1]$, in $U=(a+1)T+2aT=(3a+1)T\le2sT\le4scL$ ticks, as
$T\le c(L+1)\le2cL$. By \cref{lem:skel-transport} with $K=K_w$, $\rho=1$, $\kappa=wB+1$, $\omega=w-1$ and $W=db\ge2w-1$,
\[
  m\log_2d\le(wB+1)K_w(K_w+4)\cdot4scL+3L,
\]
which is the claim.
\end{informalproof}

\begin{proposition}[Paired distances on a skeleton class]\label{prop:skel-paired}
\notformalized
Let $k,w\ge1$ and $D\ge0$ be integers and $c>0$. Let $\sigma$ be a skeleton at scale~$w$ of a machine with $k$ work
tapes at a length $L=d^2b(d)$ with $d\ge2$ and $db(d)\ge2w-1$, such that $\mathcal X_\sigma\ne\emptyset$,
$T_\sigma\le c(L+1)$ and every member appends exactly $L$ bits; let $\Gamma$ have full depth at most~$D$ in
$\bar H_\sigma$, let $E$ be the incidence graph of $N_\Gamma$, and for $Q\subseteq[0,L)$ let $E_Q$ be its induced
subgraph on $\{L_p,R_p:p\in Q\}$ \cite[Definition~3.2]{DirectedRecovery2026}. With $B$ and $A^{(w)}_s$ as above:
\begin{enumerate}[label=(\alph*)]
\item For all integers $r\ge1$ and $\Delta\ge0$ and every $Q\subseteq[0,L)$ such that $E_Q$ has maximum degree at
most~$\Delta$,
\[
  \card{\{p\in Q:\dist_{E_Q}(L_p,R_p)\le r\}}\le\sum_{\substack{1\le s\le r\\ s\text{ odd}}}(s+1)\,\nu_r\,A^{(w)}_s\,
  \frac L{\log_2d},\qquad\nu_r=\sum_{j=0}^r\max(1,\Delta)^j .
\]
\item The set $S_\Gamma=\{p:p\in N_\Gamma(p)\}$ of self-demands has at most $A^{(w)}_1L/\log_2d$ elements.
\item For every integer $\vartheta\ge1$, the set $\mathcal H_\vartheta=\{p:\deg_E(L_p)>\vartheta\}$ has at most
$wBL/\vartheta$ elements, and $E_{[0,L)\setminus\mathcal H_\vartheta}$ has maximum degree at most $\max(wB,\vartheta)$.
\end{enumerate}
The bounds depend only on $k$, $D$, $c$, $w$, $\Delta$ and~$r$; no correctness and no bound on $\card\Gamma$ is assumed.
\end{proposition}

\begin{informalproof}
(a) This is the packing argument of the fifth paper's proof of its Theorem~3.3~\cite{DirectedRecovery2026}. For each
$p\in Q$ with $\dist_{E_Q}(L_p,R_p)\le r$ choose a shortest path from $L_p$ to~$R_p$ in~$E_Q$; it has an odd number
$s\le r$ of edges, because every edge of~$E$ joins an $L$-vertex to an $R$-vertex. A vertex~$y$ lies on at most $\nu_r$
chosen paths: their left ends are distinct and lie within distance~$r$ of~$y$ in a graph of maximum degree at
most~$\Delta$. For each odd $s\le r$ choose a maximal vertex-disjoint family of chosen paths with $s$ edges. Every
chosen path with $s$ edges meets a member of the family, and each member, having $s+1$ vertices, meets at most
$(s+1)\nu_r$ chosen paths. The members are paths in~$E$, so by \cref{lem:skel-disjoint} there are at most
$A^{(w)}_sL/\log_2d$ of them. Summing over~$s$ gives~(a).

(b) The edges $L_pR_p$, $p\in S_\Gamma$, are pairwise vertex-disjoint paths with one edge each; apply
\cref{lem:skel-disjoint} with $s=1$.

(c) Counting edges, $\sum_p\deg_E(L_p)=\sum_q\card{N_\Gamma(q)}\le wBL$ by \cref{lem:skel-cones}(a), so at most
$wBL/\vartheta$ coordinates have $\deg_E(L_p)>\vartheta$. Deleting both ports of these coordinates does not raise any
degree; the remaining left degrees are at most~$\vartheta$, and the right degrees $\deg_E(R_q)=\card{N_\Gamma(q)}$ are
at most $wB$.
\end{informalproof}

\subsection{Theorem SE}\label{sec:skel-SE}

The \emph{induced relation} of a relation~$N$ on $Q\subseteq[0,L)$ is $N_Q(q)=N(q)\cap Q$ for $q\in Q$
\cite[Section~2]{DirectedRecovery2026}.

\begin{lemma}[A uniform half-rate bound on class relations]\label{lem:skel-uniform}
\notformalized
Assume [Q9.1], an affirmative answer to Question~9.1 of the fifth paper~\cite{DirectedRecovery2026} for every fixed
degree bound; the binary alphabet suffices. Let $M$ be a machine with $k$ work tapes, let $w\ge1$ and $D\ge0$ be
integers, $c>0$ and $0<\delta<1$, and put $B=\sum_{h=0}^D(k+1)^h$ and $\vartheta=\lceil wB/\delta\rceil$. For a length
$L=d^2b(d)$ let $\mathcal N_L$ be the set of the induced relations $(N_\Gamma)_Q$, on the coordinate sets
$Q=[0,L)\setminus(\mathcal H_\vartheta\cup S_\Gamma)$, for all skeletons~$\sigma$ of~$M$ at scale~$w$ at length~$L$ such
that $\mathcal X_\sigma\ne\emptyset$, $T_\sigma\le c(L+1)$ and every member appends exactly $L$ bits, and all
$\Gamma\subseteq[0,T_\sigma)$ of full depth at most~$D$ in $\bar H_\sigma$; here $S_\Gamma$ and $\mathcal H_\vartheta$
are as in \cref{prop:skel-paired}. Then for every $\zeta>0$ there is $L_\zeta$ such that for every length
$L\ge L_\zeta$, every relation of $\mathcal N_L$, on its coordinate set~$Q$, and every binary storage code
$\mathcal C\subseteq\{0,1\}^Q$ for it, $\log_2\card{\mathcal C}\le(\frac12+\zeta)\card Q$.
\end{lemma}

\begin{informalproof}
Suppose not. Then there are $\zeta>0$ and infinitely many lengths~$L$, each with a relation $N'_L\in\mathcal N_L$ on a
coordinate set~$Q_L$, coming from some $\sigma$ and~$\Gamma$, and a binary storage code $\mathcal C_L$ for $N'_L$ with
$\log_2\card{\mathcal C_L}>(\frac12+\zeta)\card{Q_L}$. Along these lengths $d\to\infty$, so for all large ones $d\ge2$
and $db(d)\ge2w-1$, and \cref{prop:skel-paired} applies. By its parts~(b) and~(c),
$\card{Q_L}\ge(1-\delta-A^{(w)}_1/\log_2d)L$, which tends to infinity. The incidence graph of $N'_L$ is $E_{Q_L}$, which
has maximum degree at most $\max(wB,\vartheta)=\vartheta$ by part~(c), since $\vartheta\ge wB/\delta>wB$. For every
fixed $r\ge1$, part~(a) with $\Delta=\vartheta$ and $Q=Q_L$ shows that at most $O_r(L/\log_2d)=o(\card{Q_L})$
coordinates $p\in Q_L$ have $\dist_{E_{Q_L}}(L_p,R_p)\le r$; this is the paired-distance condition
\cite[Definition~2.3]{DirectedRecovery2026}. So the relations $N'_L$, the degree bound~$\vartheta$ and the codes
$\mathcal C_L$ satisfy the hypotheses of Question~9.1 with the binary alphabet, and its affirmative answer gives
$\limsup\log_2\card{\mathcal C_L}/\card{Q_L}\le\frac12$, a contradiction. The relations are loopless, since $Q_L$ avoids
$S_\Gamma$, so the argument does not depend on whether recovery relations in Question~9.1 may contain loops. This is the
uniformity argument of the fifth paper's Proposition~9.2~\cite{DirectedRecovery2026}.
\end{informalproof}

\begin{theorem}[Theorem SE: skeleton entropy]\label{thm:SE}
\notformalized
Assume [Q9.1], an affirmative answer to Question~9.1 of the fifth paper~\cite{DirectedRecovery2026} for every fixed
degree bound; the binary alphabet suffices. Let $k,V,w\ge1$ be integers, $c>0$ and $0<\eta<\frac12$. For every
machine~$M$ with $k$ work tapes there is~$L_0$ such that at no length $L=d^2b(d)\ge L_0$ does $M$ have the following
three properties:
\begin{enumerate}[label=(\roman*)]
\item on every input $X\in\{0,1\}^L$ it halts within $c(L+1)$ steps with output $X^\T$;
\item on every input, every work-tape cell is entered at most $V$ times;
\item $\mathrm H(\sigma_w(X))\le(\frac12-\eta)L$.
\end{enumerate}
More precisely, at no length $L\ge L_0$ is there a skeleton class $\mathcal X_\sigma$ of~$M$ at scale~$w$ with at least
$2^{(\frac12+\eta)L}$ elements such that $M$ halts within $c(L+1)$ steps with output $X^\T$ on every
$X\in\mathcal X_\sigma$ and has entry at most~$V$ on every work tape on some $X\in\mathcal X_\sigma$.
\end{theorem}

\begin{informalproof}
The second statement implies the first: by~\eqref{eq:skel-heavy}, (iii) gives a class with at least
$2^{(\frac12+\eta)L}$ elements, (i) holds on it, and (ii) holds on each of its members.

Put $\eps=\delta=\eta/2$. Let $\beta_w$ be the value width of the block machines of~$M$, let $D$ be given by
\cref{lem:skel-genie} for $k$, $V$, $w$, $c$ and~$\eps$, and let $B=\sum_{h=0}^D(k+1)^h$ and
$\vartheta=\lceil wB/\delta\rceil$; all depend only on $M$, $V$, $w$, $c$ and~$\eta$. Let $L_{\eta/8}$ be given by
\cref{lem:skel-uniform} for these $D$ and~$\delta$, and let $L_0\ge L_{\eta/8}$ be such that every length $L\ge L_0$ has
$d\ge2$, $db(d)\ge2w-1$ and $A^{(w)}_1/(2\log_2d)<\eta/8$.

Suppose that a class $\mathcal X_\sigma$ as described exists at a length $L\ge L_0$. Its members append exactly $L$ bits
and halt within $c(L+1)$ steps, so $T_\sigma\le c(L+1)$. By \cref{lem:skel-genie} there is a genie~$\Gamma$ of full
depth at most~$D$ in $\bar H_\sigma$ with $\beta_w\card\Gamma\le\eps L$. The revealed word $G_\Gamma$ takes at most
$2^{\eps L}$ values, so some fibre $W=\{X\in\mathcal X_\sigma:G_\Gamma(X)=g\}$ has
\[
  \log_2\card W\ge\bigl(\tfrac12+\eta\bigr)L-\eps L=\bigl(\tfrac12+\tfrac\eta2\bigr)L .
\]
By \cref{lem:skel-cones}(c), $W$ is a binary storage code for~$N_\Gamma$. Let
$Q=[0,L)\setminus(\mathcal H_\vartheta\cup S_\Gamma)$. By the fifth paper's paid
conditioning~\cite[Lemma~2.2]{DirectedRecovery2026}, which is Lean-verified there, fixing the coordinates in
$\mathcal H_\vartheta\cup S_\Gamma$ and restricting to~$Q$ turns every nonempty fibre of~$W$ into a storage code for
$(N_\Gamma)_Q\in\mathcal N_L$, which has at most $2^{(\frac12+\eta/8)\card Q}$ words by \cref{lem:skel-uniform} with
$\zeta=\eta/8$; and there are at most $2^{\card{\mathcal H_\vartheta\cup S_\Gamma}}$ such fibres. By
\cref{prop:skel-paired}(b),(c), $\card{S_\Gamma}\le A^{(w)}_1L/\log_2d$ and
$\card{\mathcal H_\vartheta}\le wBL/\vartheta\le\delta L=\eta L/2$. Hence
\[
  \log_2\card W\le\card{\mathcal H_\vartheta\cup S_\Gamma}+\bigl(\tfrac12+\tfrac\eta8\bigr)
  \bigl(L-\card{\mathcal H_\vartheta\cup S_\Gamma}\bigr)\le\bigl(\tfrac12+\tfrac\eta8\bigr)L
  +\tfrac12\card{\mathcal H_\vartheta\cup S_\Gamma}\le\bigl(\tfrac12+\tfrac{3\eta}8\bigr)L+\frac{A^{(w)}_1L}{2\log_2d},
\]
which is less than $(\frac12+\frac\eta2)L$, a contradiction.
\end{informalproof}

The hypotheses enter the proof as follows: (iii) only through the existence of one large class; (ii) only for one member
of that class, through its coarse work walks; and (i) for the time bound, for the number of appended bits and for
correctness. Given [Q9.1], Theorem~SE is a necessary condition for adaptive transposers: a linear-time transposer whose
work heads enter every cell boundedly often must, at every fixed scale~$w$ and for every $\eta>0$, have skeleton entropy
above $(\frac12-\eta)L$ at all large lengths. Its cut is determined by the skeleton alone, so no structure of the
individual traces is paid for, and motion inside blocks, however data-dependent, is invisible to it.

\paragraph{Relation to Corollary F$''$.} Theorem~SE contains the eighth paper's Corollary~F$''$
\cite[Corollary~6.4]{BoundedNesting2026}. Suppose that $M$ has the three properties of that corollary at a length~$L$,
with a constant $\eta'<\frac12$ in place of~$\eta$. Some trace~$\tau_0$ has probability at least $2^{-\eta'L}$, and it
enters every work-tape cell at most $V$ times. A trace determines the skeleton, so the at least $2^{(1-\eta')L}$ inputs
that follow~$\tau_0$ lie in one skeleton class at any scale~$w$. With $\eta=\min(\frac12-\eta',\frac14)$, this class has
at least $2^{(\frac12+\eta)L}$ elements and contains an input on which $M$ has entry at most~$V$ on every work tape,
which the second statement of Theorem~SE excludes at all large lengths.

\subsection{Bounded entry and input steering}\label{sec:skel-steering}

\begin{lemma}[Few coarse walks]\label{lem:skeleton-count}
\notformalized
\begin{enumerate}[label=(\alph*)]
\item Let $V\ge1$, $w\ge4V$ and $T\ge16V$ be integers. The coarse walks $\lfloor x/w\rfloor$ of all walks~$x$ of
length~$T$ with $x(0)=0$ that enter every cell at most $V$ times are at most $2^{m_1\log_2(ew/V)}$ in number, where
$m_1=2V(T/w+2)$.
\item Let $T\ge2$ and $1\le m\le T-1$ be integers. The maps $s\colon[0,T)\to\Z$ with $s(0)=0$, $s(t+1)-s(t)\in\{0,1\}$
for $t+1<T$, and $s(T-1)\le m$ are at most $(eT/y)^y$ in number, for every real number~$y$ with $m\le y\le T$.
\end{enumerate}
\end{lemma}

\begin{informalproof}
We use $\sum_{i=0}^{m}\binom ni\le(en/m)^m$ for integers $1\le m\le n$.

(a) By \cref{lem:skel-coarse}(c), $s=\lfloor x/w\rfloor$ changes value at most $m_0=2V((T-1)/w+2)-1$ times, each time
by~$\pm1$, at times in $[1,T)$; since $s(0)=0$, $s$ is determined by the set of these times and the signs of the
changes. Let $m_0'=\lfloor m_0\rfloor$. Then $m_0'\ge4V-1\ge3$, and $m_0'\le(T-1)/2+4V-1\le T-1$, because $w\ge4V$ and
$T\ge8V-1$. So the number of coarse walks is at most
\[
  \sum_{i=0}^{m_0'}\binom{T-1}i2^i\le2^{m_0'}\Bigl(\frac{e(T-1)}{m_0'}\Bigr)^{m_0'}\le\Bigl(\frac{2eT}{m_0'}\Bigr)^{m_0'} .
\]
The function $y\mapsto y\ln(2eT/y)$ increases for $0<y<2T$, its derivative being $\ln(2T/y)$, and
$m_0'<m_1\le T/2+4V\le3T/4$, because $w\ge4V$ and $T\ge16V$. So the count is at most $(2eT/m_1)^{m_1}\le(ew/V)^{m_1}$,
because $m_1\ge2VT/w$.

(b) Such a map is determined by the set of times $t\in[1,T)$ with $s(t)\ne s(t-1)$, which has at most $m$ elements. So
there are at most $\sum_{i\le m}\binom{T-1}i\le(e(T-1)/m)^m\le(eT/m)^m$ of them, and $y\mapsto y\ln(eT/y)$ increases on
$(0,T]$, its derivative being $\ln(T/y)$.
\end{informalproof}

A head that enters every cell at most $V$ times has a coarse walk that carries only $O(V\log(w/V)/w)$ bits per step; so
at large scales the skeleton of a machine with bounded entry is almost determined by its length, and only heads that
revisit cells unboundedly often can carry a constant fraction of the input in their coarse motion.

\begin{lemma}[Coarse entropy of bounded-entry heads]\label{lem:skel-entropy}
\notformalized
Let $M$ have $k$ work tapes, let $V\ge1$ and $w\ge4V$ be integers, $c>0$, and let $L\ge\max(2w,16V)$ be a length.
Suppose that on every input $X\in\{0,1\}^L$, $M$ appends exactly $L$ bits, halts within $c(L+1)$ steps and has entry at
most~$V$ on every work tape. Then
\[
  \mathrm H(\sigma_w(X))\le\mathrm H(\sigma^{\mathrm{in}}_w(X))+\theta_{k,V,c}(w)\,L+4kV\log_2\frac{ew}V,
\]
where $\theta_{k,V,c}(w)=\bigl(4kVc\log_2(ew/V)+\log_2(2ecw)\bigr)/w$. For fixed $k$, $V$ and~$c$,
$\theta_{k,V,c}(w)\to0$ as $w\to\infty$.
\end{lemma}

\begin{informalproof}
The skeleton $\sigma_w(X)$ is determined by the input skeleton, which fixes $T=T(X)$, together with the $k$ coarse work
walks and the output component. So $\mathrm H(\sigma_w(X))\le\mathrm H(\sigma^{\mathrm{in}}_w(X))+\log_2N$, where $N$ is
the largest number, over the values of~$T$, of possible values of these $k+1$ components given~$T$. Since $L$ bits are
appended at distinct steps, $L\le T\le c(L+1)\le2cL$; in particular $T\ge16V$. Each coarse work walk is the coarse walk
of a walk of length~$T$ that starts at~$0$ and enters every cell at most $V$ times, so by \cref{lem:skeleton-count}(a)
it takes at most $2^{m_1\log_2(ew/V)}$ values, with $m_1=2V(T/w+2)\le4VcL/w+4V$. The output component starts at~$0$,
grows by $0$ or~$1$ per step, and ends at most at $m=\lfloor L/w\rfloor$, where $1\le m\le L/w\le T-1$ because $L\ge2w$
and $w\ge4$; by \cref{lem:skeleton-count}(b) with $y=L/w$ it takes at most $(eTw/L)^{L/w}\le(2ecw)^{L/w}$ values. Hence
$\log_2N\le k(4VcL/w+4V)\log_2(ew/V)+(L/w)\log_2(2ecw)$, which is the claim. The limit holds because $(\log_2w)/w\to0$.
\end{informalproof}

The output component alone contributes at most $(L/w)\log_2(2ecw)$ bits, whatever the machine does, as long as it
appends exactly $L$ bits within $c(L+1)$ steps; it could therefore be omitted from the skeleton, at the price of passing
to a coarser scale and a smaller~$\eta$ in Theorem~SE.

\begin{corollary}[Corollary IS: input steering]\label{cor:IS}
\notformalized
Assume [Q9.1], an affirmative answer to Question~9.1 of the fifth paper~\cite{DirectedRecovery2026} for every fixed
degree bound; the binary alphabet suffices. Let $k,V,w\ge1$ be integers, $c>0$ and $0<\eta<\frac12$. For every
machine~$M$ with $k$ work tapes there is~$L_0$ such that at no length $L=d^2b(d)\ge L_0$ does $M$ have the following
three properties:
\begin{enumerate}[label=(\roman*)]
\item on every input $X\in\{0,1\}^L$ it halts within $c(L+1)$ steps with output $X^\T$;
\item on every input, every work-tape cell is entered at most $V$ times;
\item $\mathrm H(\sigma^{\mathrm{in}}_w(X))\le(\frac12-\eta)L$.
\end{enumerate}
\end{corollary}

\begin{informalproof}
Choose an integer $m\ge1$ such that $w'=mw$ satisfies $w'\ge4V$ and $\theta_{k,V,c}(w')\le\eta/2$, which is possible by
\cref{lem:skel-entropy}. Let $L_0$ be at least the length that Theorem~SE gives for $M$, $V$, $w'$, $c$ and $\eta/4$,
and such that every length $L\ge L_0$ satisfies $L\ge\max(2w',16V)$ and $4kV\log_2(ew'/V)\le\eta L/4$. Suppose that $M$
has (i)--(iii) at a length $L\ge L_0$. Since $\lfloor z/(mw)\rfloor=\lfloor\lfloor z/w\rfloor/m\rfloor$ for every
integer~$z$, the input skeleton at scale~$w'$ is a function of that at scale~$w$, so
$\mathrm H(\sigma^{\mathrm{in}}_{w'}(X))\le(\frac12-\eta)L$. By (i), (ii) and \cref{lem:skel-entropy} at scale~$w'$,
\[
  \mathrm H(\sigma_{w'}(X))\le\bigl(\tfrac12-\eta\bigr)L+\tfrac\eta2L+\tfrac\eta4L=\bigl(\tfrac12-\tfrac\eta4\bigr)L .
\]
So $M$ has the three properties of Theorem~SE at scale~$w'$ with $\eta/4$ in place of~$\eta$, which is impossible.
\end{informalproof}

Corollary~IS reduces adaptivity, for machines whose work tapes have bounded entry, to the coarse motion of one head:
given [Q9.1], such a linear-time transposer must move its input head, at every fixed scale, in a way that depends on the
data and carries about half of the input. For a construction this is a requirement; for a proof of $\neg$(T) it leaves,
among such machines, only those whose input head is steered in this way. For an oblivious machine the input skeleton is
constant, so Corollary~IS contains the eighth paper's Corollary~F$'$~\cite[Corollary~6.3]{BoundedNesting2026} and
extends it to adaptive machines.

\begin{corollary}[Corollary RV: revisits]\label{cor:RV}
\notformalized
Assume [Q9.1], an affirmative answer to Question~9.1 of the fifth paper~\cite{DirectedRecovery2026} for every fixed
degree bound; the binary alphabet suffices. For all integers $k,V\ge1$ and every $c>0$, no machine with $k$ work tapes
is, at infinitely many lengths $L=d^2b(d)$, a transposer of all $d\times d$ matrices within $c(L+1)$ steps such that on
every input every cell of every tape, the input tape included, is entered at most $V$ times.
\end{corollary}

\begin{informalproof}
Let $M$ be such a machine. Since $(\log_2w)/w\to0$, there is an integer $w\ge4V$ with $4Vc\log_2(ew/V)\le w/8$; fix it.
At a length $L\ge16V$ at which $M$ has the stated properties, every input satisfies $L\le T(X)\le c(L+1)$, and the input
skeleton at scale~$w$ is the coarse walk of a walk of length $T(X)$ that starts at~$0$ and enters every cell at most $V$
times. By \cref{lem:skeleton-count}(a), for each value of $T(X)$ it takes at most $2^{m_1\log_2(ew/V)}$ values, where
\[
  m_1=2V\bigl(T(X)/w+2\bigr)\le2V(2cL/w+2),
\]
and $T(X)$ takes at most $c(L+1)+1$ values. Hence
\[
  \mathrm H(\sigma^{\mathrm{in}}_w(X))\le\frac L8+4V\log_2\frac{ew}V+\log_2\bigl(c(L+1)+1\bigr)\le\frac L4
\]
at all large such lengths. The work tapes satisfy~(ii) of Corollary~IS, so Corollary~IS with this~$w$ and $\eta=\frac14$
excludes $M$ at all large lengths.
\end{informalproof}

Given [Q9.1], every linear-time transposer, adaptive or not, revisits cells unboundedly often: for every~$V$, at all
large lengths, some input makes some head enter some cell more than $V$ times. Corollary~RV is incomparable with the
eighth paper's Corollary~F$'$: it needs bounded entry on the input tape, which F$'$ does not, and no obliviousness,
which F$'$ needs; the extension of F$'$ to adaptive machines is Corollary~IS. Since bounded entry allows unboundedly
many reversals~\cite[Remark~6.2]{BoundedNesting2026}, Corollary~RV does not follow from reversal or pass bounds. It is
specific to transposition, which enters through the row potential of \cref{lem:skel-transport}: copying the input to the
output takes linear time and enters every cell of every tape at most once.

\begin{remark}[Padding]\label{rem:skel-padding}
(a) Padding of bounded entry before the run does not affect Corollary~IS. Let $c_0,V_0\ge1$ be integers, and let $M'$
first move its input head along a walk that starts and ends at cell~$0$, has length between $L$ and $c_0L$ and enters
every cell at most $V_0$ times, chosen in any way as a function of the input, and then run~$M$ from its initial
configuration. The padding of the eighth paper~\cite[Section~6.2]{BoundedNesting2026}, a sweep across the input with a
step out and back at every cell of the first half that holds a~$1$, followed by a return to cell~$0$, is of this kind
with $c_0=3$ and $V_0=4$. The prefix does not move the work heads, so $M'$ inherits~(ii) of Corollary~IS from~$M$, and
it satisfies~(i) with $c+c_0$ in place of~$c$. The input skeleton of~$M'$ at a scale~$w'$ is the coarse walk of the
prefix followed by the input skeleton $\sigma^{\mathrm{in}}_{w'}(X;M)$ of~$M$. So, by \cref{lem:skeleton-count}(a)
applied to each length of the prefix, for $w'\ge4V_0$ and $L\ge16V_0$,
\[
  \mathrm H(\sigma^{\mathrm{in}}_{w'}(X;M'))\le\mathrm H(\sigma^{\mathrm{in}}_{w'}(X;M))
  +2V_0\Bigl(\frac{c_0L}{w'}+2\Bigr)\log_2\frac{ew'}{V_0}+\log_2(c_0L+1).
\]
If $M$ satisfies~(iii) of Corollary~IS at a scale~$w$ with~$\eta$, then for every multiple $w'\ge4V_0$ of~$w$ with
$2V_0c_0\log_2(ew'/V_0)/w'\le\eta/2$, the machine~$M'$ satisfies~(iii) at~$w'$ with $\eta/4$ at all large lengths. So
Corollary~IS excludes~$M'$ whenever its hypotheses~(ii) and~(iii) hold for~$M$.

(b) Padding in which the input head revisits cells unboundedly often can break hypothesis~(iii) of Theorem~SE and of
Corollary~IS at every scale. Before running~$M$, a machine may copy its input to an extra work tape and then, reading
the copy, step its input head from cell~$0$ to cell~$-1$ and back at every~$1$ and pause at every~$0$. This costs $O(L)$
steps and keeps the entries on the work tapes bounded, and $\sigma^{\mathrm{in}}_w(X)$ then determines~$X$ for
every~$w$; but the input head enters the cells $0$ and~$-1$ once for each~$1$ of the input. By~(a), padding before the
run needs such revisits to have this effect at large scales. So Theorem~SE and Corollary~IS are necessary conditions
that padding of this kind meets at no cost to a construction. Corollary~RV, whose hypotheses are entry bounds and
involve no entropy, is not affected by padding: padding adds entries and removes none.
\end{remark}

\begin{remark}[A weak coding bound suffices for Corollary RV]\label{rem:skel-weak-crux}
Proof sketch only: Corollary~RV does not need rate~$\frac12$. Under its hypotheses all $k+1$ heads have bounded entry,
so \cref{lem:skel-entropy}, together with the bound on the input skeleton in the proof of Corollary~RV, gives
$\mathrm H(\sigma_w(X))\le\theta'(w)L+O(\log L)$ with $\theta'(w)=(4(k+1)Vc\log_2(ew/V)+\log_2(2ecw))/w$, and the
largest skeleton class has at least $2^{(1-\theta'(w))L-O(\log L)}$ elements. Suppose that along every sequence of
relations of incidence degree at most~$\Delta$ that satisfies the paired-distance condition, all binary storage codes
have rate at most $1-\rho(\Delta)+o(1)$. The proof of Theorem~SE, run on that class with $\delta=\eps$, then gives a
contradiction as soon as $\rho(\Delta(w,\eps))>\theta'(w)+2\eps$ for some $w\ge4V$ and $\eps\in(0,1)$, where
$\Delta(w,\eps)=\lceil wB/\eps\rceil$ is the trimming threshold and $B$ the cone bound at the depth of
\cref{lem:skel-genie} for~$\eps$; here $\log\log\Delta(w,\eps)$ is of order $k^3V^3\beta_wC/\eps$, and $\beta_w$ is of
order $kw\log\card\Sigma$ for large~$w$. Compare the eighth paper's Remark~6.5~\cite{BoundedNesting2026}.
\end{remark}
